\Chapter{21}

\Verse{1}
{
    \zA{话说史湘云说着笑着跑出来，怕黛玉赶上。 宝玉在后忙说：“绊倒了!那里就 赶上了？” 黛玉赶到门前，被宝玉叉手在门框上拦住，笑道：“饶他这一遭儿罢。”
        黛玉拉着手说道：“我要饶了云儿，再不活着。” 湘云见宝玉拦着门，料黛玉不能 出来，便立住脚，笑道：“好姐姐，饶我这遭儿罢！” 却值宝钗来在湘云身背后，
        也笑道：“我劝你们两个看宝兄弟面上，都撂开手罢。” 黛玉道：“我不依。 你们 是一气的，都来戏弄我。” 宝玉劝道：“罢呦，谁敢戏弄你？}
}
{
    \eA{But to resume our story. When Shih Hsiang-yün ran out of the room, she was
        all in a flutter lest Lin Tai-yü should catch her up; but Pao-yü, who came
        after her, readily shouted out, "You'll trip and fall. How ever could she
        come up to you?" Lin Tai-yü went in pursuit of her as far as the entrance,
        when she was impeded from making further progress by Pao-yü, who stretched
        his arms out against the posts of the door.}
}
{
}

\Verse{2}
{
    \zA{你不打趣他，他就敢 说你了？” 四人正难分解，有人来请吃饭，方往前边来。 那天已掌灯时分，王夫人、 李纨、凤姐、迎探惜姊妹等，都往贾母这边来。
        大家闲话了一回，各自归寝。 湘云 仍往黛玉房中安歇。 宝玉送他二人到房，那天已二更多了，袭人来催了几次方回。 次早，天方明时， 便披衣？
        鞋往黛玉房中来了，却不见紫鹃翠缕二人，只有他姊妹两个尚卧在衾内。 那黛玉严严密密裹着一幅杏子红绫被，安稳合目而睡。}
}
{
    \eA{"Were I to spare Yün Erh, I couldn't live!" Lin Tai-yü exclaimed, as she
        tugged at his arms. But Hsiang-yün, perceiving that Pao-yü obstructed the
        door, and surmising that Tai-yü could not come out, speedily stood still.
        "My dear cousin," she smilingly pleaded, "do let me off this time!"}
}
{
}

\Verse{3}
{
    \zA{湘云却一把青丝，拖于枕畔， 一幅桃红绸被只齐胸盖着，衬着那一弯雪白的膀子，撂在被外，上面明显着两个金 镯子。
        宝玉见了叹道：“睡觉还是不老实!回来风吹了，又嚷肩膀疼了。” 一面说， 一面轻轻的替他盖上。 黛玉早已醒了，觉得有人，就猜是宝玉，翻身一看，果然是 他。
        因说道：“这早晚就跑过来作什么？” 宝玉说道：“这还早呢!你起来瞧瞧罢。” 黛玉道：“你先出去，让我们起来。” 宝玉出至外间。 黛玉起来，叫醒湘云，二人
        都穿了衣裳。}
}
{
    \eA{But it just happened that Pao-ch'ai, who was coming along, was at the back
        of Hsiang-yün, and with a face also beaming with smiles: "I advise you
        both," she said, "to leave off out of respect for cousin Pao-yü, and have
        done." "I don't agree to that," Tai-yü rejoined; "are you people, pray, all
        of one mind to do nothing but make fun of me?" "Who ventures to make fun of
        you?" Pao-yü observed advisingly;}
}
{
}

\Verse{4}
{
    \zA{宝玉又复进来坐在镜台旁边，只见紫鹃翠缕进来伏侍梳洗。 湘云洗了 脸，翠缕便拿残水要泼，宝玉道：“站着，我就势儿洗了就完了，省了又过去费事。”
        说着，便走过来，弯着腰洗了两把。 紫鹃递过香肥皂去，宝玉道：“不用了，这盆 里就不少了。” 又洗了两把，便要手巾。 翠缕撇嘴笑道：“还是这个毛病儿。”
        宝 玉也不理他，忙忙的要青盐擦了牙，漱了口。 完毕，见湘云已梳完了头，便走过来 笑道：“好妹妹，替我梳梳呢。” 湘云道：“这可不能了。”}
}
{
    \eA{"and hadn't you made sport of her, would she have presumed to have said
        anything about you?"}
}
{
}

\Verse{5}
{
    \zA{宝玉笑道：“好妹妹， 你先时候儿怎么替我梳了呢？” 湘云道：“如今我忘了，不会梳了。” 宝玉道：“横 竖我不出门，不过打几根辫子就完了。”
        说着，又千“妹妹”万“妹妹”的央告。 湘云只得扶过他的头来梳篦。 原来宝玉在家并不戴冠，只将四围短发编成小辫，往 顶心发上归了总，编一根大辫，红绦结住。
        自发顶至辫梢，一路四颗珍珠，下面又 有金坠脚儿。 湘云一面编着，一面说道：“这珠子只三颗了，这一颗不是了。 我记 得是一样的，怎么少了一颗？”}
}
{
    \eA{While this quartet were finding it an arduous task to understand one
        another, a servant came to invite them to have their repast, and they
        eventually crossed over to the front side, and as it was already time for
        the lamps to be lit, madame Wang, widow Li Wan, lady Feng, Ying Ch'un, T'an
        Ch'un, Hsi Ch'un and the other cousins, adjourned in a body to dowager lady
        Chia's apartments on this side, where the whole company spent a while in a
        chat on irrelevant topics, after which they each returned to their rooms and
        retired to bed.}
}
{
}

\Verse{6}
{
    \zA{宝玉道：“丢了一颗。” 湘云道：“必定是外头去， 掉下来，叫人拣了去了。 倒便宜了拣的了。” 黛玉旁边冷笑道：“也不知是真丢，
        也不知是给了人镶什么戴去了呢！” 宝玉不答，因镜台两边都是妆奁等物，顺手拿 起来赏玩，不觉拈起了一盒子胭脂，意欲往口边送，又怕湘云说。 正犹豫间，湘云
        在身后伸过手来，“拍”的一下将胭脂从他手中打落，说道：“不长进的毛病儿! 多早晚才改呢？”
        一语未了，只见袭人进来，见这光景，知是梳洗过了，只得回来自己梳洗。}
}
{
    \eA{Hsiang-yün, as of old, betook herself to Tai-yü's quarters to rest, and
        Pao-yü escorted them both into their apartment, and it was after the hour
        had already past the second watch, and Hsi Jen had come and pressed him
        several times, that he at length returned to his own bedroom and went to
        sleep.}
}
{
}

\Verse{7}
{
    \zA{忽 见宝钗走来，因问：“宝兄弟那里去了？” 袭人冷笑道：“‘宝兄弟’那里还有在 家的工夫！” 宝钗听说，心中明白。
        袭人又叹道：“姐妹们和气，也有个分寸儿， 也没个黑家白日闹的。 凭人怎么劝，都是耳旁风。” 宝钗听了，心中暗忖道：“倒
        别看错了这个丫头，听他说话，倒有些识见。” 宝钗便在炕上坐了，慢慢的闲言中， 套问他年纪家乡等语，留神窥察其言语志量，深可敬爱。
        一时宝玉来了，宝钗方出去。}
}
{
    \eA{The next morning, as soon as it was daylight, he threw his clothes over him,
        put on his low shoes and came over into Tai-yü's room, where he however saw
        nothing of the two girls Tzu Chüan and Ts'ui Lu, as there was no one else
        here in there besides his two cousins, still reclining under the coverlets.
        Tai-yü was closely wrapped in a quilt of almond-red silk, and lying quietly,
        with closed eyes fast asleep;}
}
{
}

\Verse{8}
{
    \zA{宝玉便问袭人道：“怎么宝姐姐和你说的这么热 闹，见我进来就跑了？” 问一声不答。 再问时，袭人方道：“你问我吗？ 我不知道 你们的原故。”
        宝玉听了这话，见他脸上气色非往日可比，便笑道：“怎么又动了 气了呢？” 袭人冷笑道：“我那里敢动气呢？ 只是你从今别进这屋子了，横竖有人
        伏侍你，再不必来支使我。 我仍旧还伏侍老太太去。” 一面说，一面便在炕上合眼 倒下。 宝玉见了这般景况，深为骇异，禁不住赶来央告。
        那袭人只管合着眼不理。}
}
{
    \eA{while Shih Hsiang-yün, with her handful of shiny hair draggling along the
        edge of the pillow, was covered only up to the chest, and outside the
        coverlet rested her curved snow-white arm, with the gold bracelets, which
        she had on. At the sight of her, Pao-yü heaved a sigh. "Even when asleep,"
        he soliloquised, "she can't be quiet!}
}
{
}

\Verse{9}
{
    \zA{宝玉没了主意，因见麝月进来，便问道：“你姐姐怎么了？” 麝月道：“我知道么？ 问你自己就明白了。”
        宝玉听说，呆了一回，自觉无趣，便起身嗳道：“不理我罢! 我也睡去。” 说着，便起身下炕，到自己床上睡下。 袭人听他半日无动静，微微的打？
        ，料他睡着，便起来拿了一领斗篷来替他盖 上。 只听“唿”的一声，宝玉便掀过去，仍合着眼装睡。 袭人明知其意，便点头冷
        笑道：“你也不用生气，从今儿起，我也只当是个哑吧，再不说你一声儿了好不好？”}
}
{
    \eA{but by and by, when the wind will have blown on her, she'll again shout that
        her shoulder is sore!"}
}
{
}

\Verse{10}
{
    \zA{宝玉禁不住起身问道：“我又怎么了？ 你又劝我？ 你劝也罢了，刚才又没劝，我一进 来，你就不理我，赌气睡了，我还摸不着是为什么。
        这会子你又说我恼了!我何尝 听见你劝我的是什么话呢？” 袭人道：“你心里还不明白？ 还等我说呢！”
        正闹着，贾母遣人来叫他吃饭，方往前边来胡乱吃了一碗，仍回自己房中。 只 见袭人睡在外头炕上，麝月在旁抹牌。 宝玉素知他两个亲厚，并连麝月也不理，揭
        起软帘自往里间来。 麝月只得跟进来。}
}
{
    \eA{With these words, he gently covered her, but Lin Tai-yü had already awoke
        out of her sleep, and becoming aware that there was some one about, she
        promptly concluded that it must, for a certainty, be Pao-yü, and turning
        herself accordingly round, and discovering at a glance that the truth was
        not beyond her conjectures, she observed: "What have you run over to do at
        this early hour?"}
}
{
}

\Verse{11}
{
    \zA{宝玉便推他出去说：“不敢惊动。” 麝月便 笑着出来，叫了两个小丫头进去。 宝玉拿了本书，歪着看了半天，因要茶，抬头见
        两个小丫头在地下站着，那个大两岁清秀些的，宝玉问他道：“你不是叫什么‘香’ 吗？” 那丫头答道：“叫蕙香。” 宝玉又问：“是谁起的名字？”
        蕙香道：“我原 叫芸香，是花大姐姐改的。” 宝玉道：“正经叫‘晦气’也罢了，又‘蕙香’咧! 你姐儿几个？” 蕙香道：“四个。” 宝玉道：“你第几个？”
        蕙香道：“第四。”}
}
{
    \eA{to which question Pao-yü replied: "Do you call this early? but get up and
        see for yourself!" "First quit the room," Tai-yü suggested, "and let us get
        up!" Pao-yü thereupon made his exit into the ante-chamber, and Tai-yü jumped
        out of bed, and awoke Hsiang-yün. When both of them had put on their
        clothes, Pao-yü re-entered and took a seat by the side of the toilet table;}
}
{
}

\Verse{12}
{
    \zA{宝玉道：“明日就叫‘四儿’，不必什么‘蕙’香‘兰’气的。 那一个配比这些花 儿？ 没的玷辱了好名好姓的！” 一面说，一面叫他倒了茶来。
        袭人和麝月在外间听 了半日，只管悄悄的抿着嘴儿笑。 这一日，宝玉也不出房，自己闷闷的，只不过拿书解闷，或弄笔墨，也不使唤 众人，只叫四儿答应。
        谁知这四儿是个乖巧不过的丫头，见宝玉用他，他就变尽方 法儿笼络宝玉。 至晚饭后，宝玉因吃了两杯酒，眼饧耳热之馀，若往日则有袭人等 大家嘻笑有兴；}
}
{
    \eA{whence he beheld Tzu-chüan and Hsüeh Yen walk in and wait upon them, as they
        dressed their hair and performed their ablutions. Hsiang-yün had done
        washing her face, and Ts'üi Lü at once took the remaining water and was
        about to throw it away, when Pao-yü interposed, saying: "Wait, I'll avail
        myself of this opportunity to wash too and finish with it, and thus save
        myself the trouble of having again to go over!"}
}
{
}

\Verse{13}
{
    \zA{今日却冷清清的，一人对灯，好没兴趣。 待要赶了他们去，又怕他 们得了意，以后越来劝了； 若拿出作上人的光景镇唬他们，似乎又太无情了。 说不
        得横着心：“只当他们死了，横竖自家也要过的。” 如此一想，却倒毫无牵挂，反 能怡然自悦。 因命四儿剪烛烹茶，自己看了一回《南华经》，至外篇《？
        箧》一则， 其文曰： 故绝圣弃智，大盗乃止；？ 玉毁珠，小盗不起。 焚符破玺，而民朴鄙； 剖斗折 衡，而民不争； 殚残天下之圣法，而民始可与论议。}
}
{
    \eA{Speaking the while, he hastily came forward, and bending his waist, he
        washed his face twice with two handfuls of water, and when Tzu Chüan went
        over to give him the scented soap, Pao-yü added: "In this basin, there's a
        good deal of it, and there's no need of rubbing any more!" He then washed
        his face with two more handfuls, and forthwith asked for a towel, and Ts'üi
        Lü exclaimed: "What!}
}
{
}

\Verse{14}
{
    \zA{擢乱六律，铄绝竽瑟，塞瞽旷 之耳，而天下始人含其聪矣； 灭文章，散五彩，胶离朱之目，而天下始人含其明矣； 毁绝钩绳，而弃规矩，？
        工垂之指，而天下始人含其巧矣。 看至此，意趣洋洋，趁着酒兴，不禁提笔续曰： 焚花散麝，而闺阁始人含其劝矣； 戕宝钗之仙姿，灰黛玉之灵窍，丧灭情意，
        而闺阁之美恶始相类矣。 彼含其劝，则无参商之虞矣； 戕其仙姿，无恋爱之心矣； 灰其灵窍，无才思之情矣。}
}
{
    \eA{have you still got this failing? when will you turn a new leaf?" But Pao-yü
        paid not so much as any heed to her, and there and then called for some
        salt, with which he rubbed his teeth, and rinsed his mouth. When he had
        done, he perceived that Hsiang-yün had already finished combing her hair,
        and speedily coming up to her, he put on a smile, and said: "My dear cousin,
        comb my hair for me!" "This can't be done!"}
}
{
}

\Verse{15}
{
    \zA{彼钗、玉、花、麝者，皆张其罗而邃其穴，所以迷惑缠 陷天下者也。 续毕，掷笔就寝。 头刚着枕，便忽然睡去，一夜竟不知所之。
        直至天明方醒，翻身看时，只见袭人和衣睡在衾上。 宝玉将昨日的事，已付之 度外，便推他说道：“起来好生睡，看冻着。” 原来袭人见他无明无夜和姐妹们鬼
        混，若真劝他，料不能改，故用柔情以警之，料他不过半日片刻，仍旧好了； 不想 宝玉竟不回转，自己反不得主意，直一夜没好生睡。}
}
{
    \eA{Hsiang-yün objected. "My dear cousin," Pao-yü continued smirkingly, "how is
        it that you combed it for me in former times?" "I've forgotten now how to
        comb it!" Hsiang-yün replied. "I'm not, after all, going out of doors,"
        Pao-yü observed, "nor will I wear a hat or frontlet, so that all that need
        be done is to plait a few queues, that's all!"}
}
{
}

\Verse{16}
{
    \zA{今忽见宝玉如此，料是他心意 回转，便索性不理他。 宝玉见他不应，便伸手替他解衣，刚解开钮子，被袭人将手 推开，又自扣了。
        宝玉无法，只得拉他的手笑道：“你到底怎么了？” 连问几声， 袭人睁眼说道：“我也不怎么着。 你睡醒了，快过那边梳洗去。 再迟了，就赶不上 了。”
        宝玉道：“我过那里去？” 袭人冷笑道：“你问我，我知道吗？ 你爱过那里 去就过那里去。 从今咱们两个人撂开手，省的鸡生鹅斗，叫别人笑话。}
}
{
    \eA{Saying this, he went on to appeal to her in a thousand and one endearing
        terms, so that Hsiang-yün had no alternative, but to draw his head nearer to
        her and to comb one queue after another, and as when he stayed at home he
        wore no hat, nor had, in fact, any tufted horns, she merely took the short
        surrounding hair from all four sides, and twisting it into small tufts, she
        collected it together over the hair on the crown of the head, and plaited a
        large queue, binding it fast with red ribbon;}
}
{
}

\Verse{17}
{
    \zA{横竖那边腻 了过来，这边又有什么‘四儿’‘五儿’伏侍你。 我们这起东西，可是‘白玷辱了 好名好姓’的！” 宝玉笑道：“你今儿还记着呢？”
        袭人道：“一百年还记着呢。 比不得你，拿着我的话当耳旁风，夜里说了，早起就忘了。” 宝玉见他娇嗔满面，
        情不可禁，便向枕边拿起一根玉簪来，一跌两段，说道：“我再不听你说，就和这 簪子一样！” 袭人忙的拾了簪子，说道：“大早起，这是何苦来？ 听不听在你，也
        不值的这么着呀。”}
}
{
    \eA{while from the root of the hair to the end of the queue, were four pearls in
        a row, below which, in the way of a tip, was suspended a golden pendant. "Of
        these pearls there are only three," Hsiang-yün remarked as she went on
        plaiting; "this isn't one like them; I remember these were all of one kind,
        and how is it that there's one short?" "I've lost one," Pao-yü rejoined.}
}
{
}

\Verse{18}
{
    \zA{宝玉道：“你那里知道我心里的急呢？” 袭人笑道：“你也知 道着急么？ 你可知道我心里是怎么着？ 快洗脸去罢。” 说着，二人方起来梳洗。
        宝玉往上房去后，谁知黛玉走来，见宝玉不在房中，因翻弄案上书看。 可巧便 翻出昨儿的《庄子》来，看见宝玉所续之处，不觉又气又笑，不禁也提笔续了一绝 云：
        无端弄笔是何人？ 剿袭《南华》庄子文。 不悔自家无见识，却将丑语诋他人! 题毕，也往上房来见贾母，后往王夫人处来。}
}
{
    \eA{"It must have dropped," Hsiang-yün added, "when you went out of doors, and
        been picked up by some one when you were off your guard; and he's now,
        instead of you, the richer for it." "One can neither tell whether it has
        been really lost," Tai-yü, who stood by, interposed, smiling the while
        sarcastically; "nor could one say whether it hasn't been given away to some
        one to be mounted in some trinket or other and worn!"}
}
{
}

\Verse{19}
{
    \zA{谁知凤姐之女大姐儿病了，正乱着请大夫诊脉。 大夫说：“替太太奶奶们道喜： 姐儿发热是见喜了，并非别症。 “王夫人凤姐听了，忙遣人问：“可好不好？” 大
        夫回道：“症虽险，却顺，倒还不妨。 预备桑虫、猪尾要紧。” 凤姐听了，登时忙 将起来：一面打扫房屋，供奉“痘疹娘娘”； 一面传与家人忌煎炒等物；
        一面命平 儿打点铺盖衣服与贾琏隔房； 一面又拿大红尺头给奶子丫头亲近人等裁衣裳。}
}
{
    \eA{Pao-yü made no reply; but set to work, seeing that the two sides of the
        dressing table were all full of toilet boxes and other such articles, taking
        up those that came under his hand and examining them. Grasping unawares a
        box of cosmetic, which was within his reach, he would have liked to have
        brought it to his lips, but he feared again lest Hsiang-yün should chide
        him.}
}
{
}

\Verse{20}
{
    \zA{外面 打扫净室，款留两位医生，轮流斟酌诊脉下药，十二日不放家去。 贾琏只得搬出外 书房来安歇。 凤姐和平儿都跟王夫人日日供奉“娘娘”。
        那贾琏只离了凤姐，便要寻事，独寝了两夜十分难熬，只得暂将小厮内清俊的 选来出火。 不想荣国府内有一个极不成材破烂酒头厨子名叫多官儿，因他懦弱无
        能，人都叫他作“多浑虫”。 二年前他父亲给他娶了个媳妇，今年才二十岁，也有 几分人材，又兼生性轻薄，最喜拈花惹草。}
}
{
    \eA{While he was hesitating whether to do so or not, Hsiang-yün, from behind,
        stretched forth her arm and gave him a smack, which sent the cosmetic flying
        from his hand, as she cried out: "You good-for-nothing! when will you mend
        those weaknesses of yours!"}
}
{
}

\Verse{21}
{
    \zA{多浑虫又不理论，只有酒有肉有钱，就 诸事不管了，所以宁荣二府之人都得入手。 因这媳妇妖调异常，轻狂无比，众人都 叫他“多姑娘儿”。
        如今贾琏在外熬煎，往日也见过这媳妇，垂涎久了，只是内惧 娇妻，外惧娈童，不曾得手。 那多姑娘儿也久有意于贾琏，只恨没空儿； 今闻贾琏
        挪在外书房来，他便没事也要走三四趟，招惹的贾琏似饥鼠一般。 少不得和心腹小 厮计议，许以金帛，焉有不允之理，况都和这媳妇子是旧交，一说便成。}
}
{
    \eA{But hardly had she had time to complete this remark, when she caught sight
        of Hsi Jen walk in, who upon perceiving this state of things, became aware
        that he was already combed and washed, and she felt constrained to go back
        and attend to her own coiffure and ablutions. But suddenly, she saw
        Pao-ch'ai come in and inquire: "Where's cousin Pao-yü gone?"}
}
{
}

\Verse{22}
{
    \zA{是夜多浑 虫醉倒在炕，二鼓人定，贾琏便溜进来相会。 一见面早已神魂失据，也不及情谈款
        叙，便宽衣动作起来，谁知这媳妇子有天生的奇趣，一经男子挨身，便觉遍体筋骨 瘫软，使男子如卧绵上，更兼淫态浪言，压倒娼妓。 贾琏此时恨不得化在他身上。
        那媳妇子故作浪语，在下说道：“你们姐儿出花儿，供着娘娘，你也该忌两日，倒 为我腌？ 了身子，快离了我这里罢。” 贾琏一面大动，一面喘吁吁答道：“你就是
        ‘娘娘’!那里还管什么‘娘娘’呢！”}
}
{
    \eA{"Do you mean to say," Hsi Jen insinuated with a sardonic smile, "that your
        cousin Pao-yü has leisure to stay at home?" When Pao-ch'ai heard these
        words, she inwardly comprehended her meaning, and when she further heard Hsi
        Jen remark with a sigh: "Cousins may well be on intimate terms, but they
        should also observe some sort of propriety; and they shouldn't night and day
        romp together;}
}
{
}

\Verse{23}
{
    \zA{那媳妇子越浪起来，贾琏亦丑态毕露。 一 时事毕，不免盟山誓海，难舍难分。 自此后，遂成相契。
        一日，大姐毒尽癍回，十二日后送了“娘娘”，合家祭天祀祖，还愿焚香，庆 贺放赏已毕，贾琏仍复搬进卧室。 见了凤姐，正是俗语云：“新婚不如远别。” 是
        夜更有无限恩爱，自不必说。 次日早起，凤姐往上屋里去后，平儿收拾外边拿进来 的衣服铺盖，不承望枕套中抖出一绺青丝来。}
}
{
    \eA{and no matter how people may tender advice it's all like so much wind
        blowing past the ears." Pao-ch'ai began, at these remarks, to cogitate
        within her mind: "May I not, possibly, have been mistaken in my estimation
        of this girl; for to listen to her words, she would really seem to have a
        certain amount of \_savoir faire\_!"}
}
{
}

\Verse{24}
{
    \zA{平儿会意，忙藏在袖内，便走到这边 房里，拿出头发来，向贾琏笑道：“这是什么东西？” 贾琏一见，连忙上来要抢。
        平儿就跑，被贾琏一把揪住，按在炕上，从手中来夺。 平儿笑道：“你这个没良心 的，我好意瞒着他来问你，你倒赌利害!等我回来告诉了，看你怎么着？” 贾琏听
        说，忙陪笑央求道：“好人，你赏我罢!我再不敢利害了。” 一语未了，忽听凤姐 声音。 贾琏此时松了不是抢又不是，只叫：“好人，别叫他知道！”}
}
{
    \eA{Pao-ch'ai thereupon took a seat on the stove-couch, and quietly, in the
        course of their conversation on one thing and another, she managed to
        ascertain her age, her native village and other such particulars, and then
        setting her mind diligently to put, on the sly, her conversation and mental
        capacity to the test, she discovered how deeply worthy she was to be
        respected and loved.}
}
{
}

\Verse{25}
{
    \zA{平儿才起身， 凤姐已走进来，叫平儿：“快开匣子，替太太找样子。” 平儿忙答应了，找时，凤
        姐见了贾琏，忽然想起来，便问平儿：“前日拿出去的东西，都收进来了没有？” 平儿道：“收进来了。” 凤姐道：“少什么不少？” 平儿道：“细细查了，没少一
        件儿。” 凤姐又道：“可多什么？” 平儿笑道：“不少就罢了，那里还有多出来的 分儿？” 凤姐又笑道：“这十几天，难保干净，或者有相好的丢下什么戒指儿、汗
        巾儿，也未可定。”}
}
{
    \eA{But in a while Pao-yü arrived, and Pao-ch'ai at once quitted the apartment.
        "How is it," Pao-yü at once inquired, "that cousin Pao-ch'ai was chatting
        along with you so lustily, and that as soon as she saw me enter, she
        promptly ran away?" Hsi Jen did not make any reply to his first question,
        and it was only when he had repeated it that Hsi Jen remarked: "Do you ask
        me? How can I know what goes on between you two?"}
}
{
}

\Verse{26}
{
    \zA{一席话，说的贾琏脸都黄了，在凤姐身背后，只望着平儿杀鸡 儿抹脖子的使眼色儿，求他遮盖。 平儿只装看不见，因笑道：“怎么我的心就和奶
        奶一样!我就怕有原故，留神搜了一搜，竟一点破绽儿都没有。 奶奶不信，亲自搜 搜。” 凤姐笑道：“傻丫头!他就有这些东西，肯叫咱们搜着？” 说着，拿了样子
        出去了。 平儿指着鼻子，摇着头儿，笑道：“这件事你该怎么谢我呢？”}
}
{
    \eA{When Pao-yü heard these words, and he noticed that the look on her face was
        so unlike that of former days, he lost no time in putting on a smile and
        asking: "Why is it that you too are angry in real earnest?" "How could I
        presume to get angry!" Hsi Jen rejoined smiling indifferently; "but you
        mustn't, from this day forth, put your foot into this room!}
}
{
}

\Verse{27}
{
    \zA{喜的贾琏眉开 眼笑，跑过来搂着，“心肝乖乖儿肉”的便乱叫起来，平儿手里拿着头发，笑道： “这是一辈子的把柄儿。 好便罢，不好咱们就抖出来。”
        贾琏笑着央告道：“你好 生收着罢，千万可别叫他知道。” 嘴里说着，瞅他不堤防，一把就抢过来，笑道： “你拿着到底不好，不如我烧了就完了事了。”
        一面说，一面掖在靴掖子内。 平儿 咬牙道：“没良心的，‘过了河儿就拆桥’，明儿还想我替你撒谎呢！” 贾琏见他 娇俏动情，便搂着求欢。}
}
{
    \eA{and as you have anyhow people to wait on you, you shouldn't come again to
        make use of my services, for I mean to go and attend to our old mistress, as
        in days of old." With this remark still on her lips, she lay herself down on
        the stove-couch and closed her eyes. When Pao-yü perceived the state of mind
        she was in, he felt deeply surprised and could not refrain from coming
        forward and trying to cheer her up.}
}
{
}

\Verse{28}
{
    \zA{平儿夺手跑出来，急的贾琏弯着腰恨道：“死促狭小娼妇 儿!一定浪上人的火来，他又跑了。” 平儿在窗外笑道：“我浪我的，谁叫你动火？
        难道图你舒服，叫他知道了，又不待见我呀！” 贾琏道：“你不用怕他!等我性子 上来，把这醋罐子打个稀烂，他才认的我呢!他防我像防贼的似的，只许他和男人
        说话，不许我和女人说话。 我和女人说话，略近些，他就疑惑，他不论小叔子、侄 儿、大的、小的，说说笑笑，就都使得了。 以后我也不许他见人！”}
}
{
    \eA{But Hsi Jen kept her eyes closed and paid no heed to him, so that Pao-yü was
        quite at a loss how to act. But espying She Yüeh enter the room, he said
        with alacrity: "What's up with your sister?" "Do I know?" answered She Yüeh,
        "examine your own self and you'll readily know!" After these words had been
        heard by Pao-yü, he gazed vacantly for some time, feeling the while very
        unhappy;}
}
{
}

\Verse{29}
{
    \zA{平儿道：“他 防你使得，你醋他使不得。 他不笼络着人，怎么使唤呢？ 你行动就是坏心，连我也 不放心，别说他呀。”
        贾琏道：“哦，也罢了么，都是你们行的是，我行动儿就存 坏心。 多早晚才叫你们都死在我手里呢！”
        正说着，凤姐走进院来，因见平儿在窗外，便问道：“要说话，怎么不在屋里 说，又跑出来隔着窗户闹，这是什么意思？” 贾琏在内接口道：“你可问他么，倒
        像屋里有老虎吃他呢。” 平儿道：“屋里一个人没有，我在他跟前作什么？”}
}
{
    \eA{but raising himself impetuously: "Well!" he exclaimed, "if you don't notice
        me, all right, I too will go to sleep," and as he spoke he got up, and,
        descending from the couch, he betook himself to his own bed and went to
        sleep.}
}
{
}

\Verse{30}
{
    \zA{凤姐 笑道：“没人才便宜呢。” 平儿听说，便道：“这话是说我么？” 凤姐便笑道：“不 说你说谁？” 平儿道：“别叫我说出好话来了！”
        说着也不打帘子，赌气往那边去 了。 凤姐自己掀帘进来，说道：“平儿丫头疯魔了，这蹄子认真要降伏起我来了! 仔细你的皮。”
        贾琏听了，倒在炕上，拍手笑道：“我竟不知平儿这么利害，从此 倒服了他了。” 凤姐道：“都是你兴的他，我只和你算账就完了。” 贾琏听了啐道：
        “你们两个人不睦，又拿我来垫喘儿了。}
}
{
    \eA{Hsi Jen noticing that he had not budged for ever so long, and that he
        faintly snored, presumed that he must have fallen fast asleep, so she
        speedily rose to her feet, and, taking a wrapper, came over and covered him.
        But a sound of "hu" reached her ear, as Pao-yü promptly threw it off and
        once again closed his eyes and feigned sleep.}
}
{
}

\Verse{31}
{
    \zA{我躲开你们就完了。” 凤姐道：“我看你 躲到那里去？” 贾琏道：“我自然有去处。” 说着就走，凤姐道：“你别走，我还 有话和你说呢。”
        不知何事，且听下回分解。 (本章完)}
}
{
    \eA{Hsi Jen distinctly grasped his idea and, forthwith nodding her head, she
        smiled coldly. "You really needn't lose your temper! but from this time
        forth, I'll become mute, and not say one word to you; and what if I do?"
        Pao-yü could not restrain himself from rising. "What have I been up to
        again," he asked, "that you're once more at me with your advice? As far as
        your advice goes, it's all well and good;}
}
{
}

\Verse{32}
{
    \zA{}
}
{
    \eA{but just now without one word of counsel, you paid no heed to me when I came
        in, but, flying into a huff, you went to sleep. Nor could I make out what it
        was all about, and now here you are again maintaining that I'm angry. But
        when did I hear you, pray, give me a word of advice of any kind?" "Doesn't
        your mind yet see for itself?" Hsi Jen replied; "and do you still expect me
        to tell you?" While they were disputing, dowager lady Chia sent a servant to
        call him to his repast, and he thereupon crossed over to the front; but
        after he had hurriedly swallowed a few bowls of rice, he returned to his own
        apartment, where he discovered Hsi Jen reclining on the outer stove-couch,
        while She Yüeh was playing with the dominoes by her side.}
}
{
}

\Verse{33}
{
    \zA{}
}
{
    \eA{Pao-yü had been ever aware of the intimacy which existed between She Yüeh
        and Hsi Jen, so that paying not the slightest notice to even She Yüeh, he
        raised the soft portiere and straightway walked all alone into the inner
        apartment. She Yüeh felt constrained to follow him in, but Pao-yü at once
        pushed her out, saying: "I don't venture to disturb you two;" so that She
        Yüeh had no alternative but to leave the room with a smiling countenance,
        and to bid two young waiting-maids go in. Pao-yü took hold of a book and
        read for a considerable time in a reclining position; but upon raising his
        head to ask for some tea, he caught sight of a couple of waiting-maids,
        standing below; the one of whom, slightly older than the other, was
        exceedingly winsome. "What's your name?"}
}
{
}

\Verse{34}
{
    \zA{}
}
{
    \eA{Pao-yü eagerly inquired. "I'm called Hui Hsiang, (orchid fragrance)," that
        waiting-maid rejoined simperingly. "Who gave you this name?" Pao-yü went on
        to ask. "I went originally under the name of Yün Hsiang (Gum Sandarac),"
        added Hui Hsiang, "but Miss Hua it was who changed it." "You should really
        be called Hui Ch'i, (latent fragrance), that would be proper; and why such
        stuff as Hui Hsiang, (orchid fragrance)?" "How many sisters have you got?"
        he further went on to ask of her. "Four," replied Hui Hsiang. "Which of them
        are you?" Pao-yü asked. "The fourth," answered Hui Hsiang. "By and by you
        must be called Ssu Erh, (fourth child)," Pao-yü suggested, "for there's no
        need for any such nonsense as Hui Hsiang (orchid fragrance) or Lan Ch'i
        (epidendrum perfume.)}
}
{
}

\Verse{35}
{
    \zA{}
}
{
    \eA{Which single girl deserves to be compared to all these flowers, without
        profaning pretty names and fine surnames!" As he uttered these words, he
        bade her give him some tea, which he drank; while Hsi Jen and She Yüeh, who
        were in the outer apartment, had been listening for a long time and laughing
        with compressed lips. Pao-yü did not, on this day, so much as put his foot
        outside the door of his room, but sat all alone sad and dejected, simply
        taking up his books, in order to dispel his melancholy fit, or diverting
        himself with his writing materials; while he did not even avail himself of
        the services of any of the family servants, but simply bade Ssu Erh answer
        his calls.}
}
{
}

\Verse{36}
{
    \zA{}
}
{
    \eA{This Ssu Erh was, who would have thought it, a girl gifted with matchless
        artfulness, and perceiving that Pao-yü had requisitioned her services, she
        speedily began to devise extreme ways and means to inveigle him. When
        evening came, and dinner was over, Pao-yü's eyes were scorching hot and his
        ears burning from the effects of two cups of wine that he had taken. Had it
        been in past days, he would have now had Hsi Jen and her companions with
        him, and with all their good cheer and laughter, he would have been enjoying
        himself. But here was he, on this occasion, dull and forlorn, a solitary
        being, gazing at the lamp with an absolute lack of pleasure.}
}
{
}

\Verse{37}
{
    \zA{}
}
{
    \eA{By and by he felt a certain wish to go after them, but dreading that if they
        carried their point, they would, in the future, come and tender advice still
        more immoderate, and that, were he to put on the airs of a superior to
        intimidate them, he would appear to be too deeply devoid of all feeling, he
        therefore, needless to say, thwarted the wish of his heart, and treated them
        just as if they were dead. And as anyway he was constrained also to live,
        alone though he was, he readily looked upon them, for the time being as
        departed, and did not worry his mind in the least on their account. On the
        contrary, he was able to feel happy and contented with his own society.}
}
{
}

\Verse{38}
{
    \zA{}
}
{
    \eA{Hence it was that bidding Ssu Erh trim the candles and brew the tea, he
        himself perused for a time the "Nan Hua Ching," and upon reaching the
        precept: "On thieves," given on some additional pages, the burden of which
        was: "Therefore by exterminating intuitive wisdom, and by discarding
        knowledge, highway robbers will cease to exist, and by taking off the jade
        and by putting away the pearls, pilferers will not spring to existence; by
        burning the slips and by breaking up the seals, by smashing the measures,
        and snapping the scales, the result will be that the people will not
        wrangle; by abrogating, to the utmost degree, wise rules under the heavens,
        the people will, at length, be able to take part in deliberation.}
}
{
}

\Verse{39}
{
    \zA{}
}
{
    \eA{By putting to confusion the musical scale, and destroying fifes and lutes,
        by deafening the ears of the blind Kuang, then, at last, will the human race
        in the world constrain his sense of hearing. By extinguishing literary
        compositions, by dispersing the five colours and by sticking the eyes of Li
        Chu, then, at length, mankind under the whole sky, will restrain the
        perception of his eyes. By destroying and eliminating the hooks and lines,
        by discarding the compasses and squares, and by amputating Kung Chui's
        fingers, the human race will ultimately succeed in constraining his
        ingenuity,"--his high spirits, on perusal of this passage, were so exultant
        that taking advantage of the exuberance caused by the wine, he picked up his
        pen, for he could not repress himself, and continued the text in this wise:
        "By burning the flower, (Hua-Hsi Jen) and dispersing the musk, (She Yüeh),
        the consequence will be that the inmates of the inner chambers will,
        eventually, keep advice to themselves.}
}
{
}

\Verse{40}
{
    \zA{}
}
{
    \eA{By obliterating Pao-ch'ai's supernatural beauty, by reducing to ashes
        Tai-yü's spiritual perception, and by destroying and extinguishing my
        affectionate preferences, the beautiful in the inner chambers as well as the
        plain will then, at length, be put on the same footing. And as they will
        keep advice to themselves, there will be no fear of any disagreement. By
        obliterating her supernatural beauty, I shall then have no incentive for any
        violent affection; by dissolving her spiritual perception, I will have no
        feelings with which to foster the memory of her talents. The hair-pin, jade,
        flower and musk (Pao-ch'ai, Tai-yü, Hsi Jen and She Yüeh) do each and all
        spread out their snares and dig mines, and thus succeed in inveigling and
        entrapping every one in the world."}
}
{
}

\Verse{41}
{
    \zA{}
}
{
    \eA{At the conclusion of this annex, he flung the pen away, and lay himself down
        to sleep. His head had barely reached the pillow before he at once fell fast
        asleep, remaining the whole night long perfectly unconscious of everything
        straight up to the break of day, when upon waking and turning himself round,
        he, at a glance, caught sight of no one else than Hsi Jen, sleeping in her
        clothes over the coverlet. Pao-yü had already banished from his mind every
        thought of what had transpired the previous day, so that forthwith giving
        Hsi Jen a push: "Get up!" he said, "and be careful where you sleep, as you
        may catch cold." The fact is that Hsi Jen was aware that he was, without
        regard to day or night, ever up to mischief with his female cousins;}
}
{
}

\Verse{42}
{
    \zA{}
}
{
    \eA{but presuming that if she earnestly called him to account, he would not mend
        his ways, she had, for this reason, had recourse to tender language to
        exhort him, in the hope that, in a short while, he would come round again to
        his better self. But against all her expectations Pao-yü had, after the
        lapse of a whole day and night, not changed the least in his manner, and as
        she really was in her heart quite at a loss what to do, she failed to find
        throughout the whole night any proper sleep.}
}
{
}

\Verse{43}
{
    \zA{}
}
{
    \eA{But when on this day, she unexpectedly perceived Pao-yü in this mood, she
        flattered herself that he had made up his mind to effect a change, and
        readily thought it best not to notice him. Pao-yü, seeing that she made no
        reply, forthwith stretched out his hand and undid her jacket; but he had
        just unclasped the button, when his arm was pushed away by Hsi Jen, who
        again made it fast herself. Pao-yü was so much at his wit's ends that he had
        no alternative but to take her hand and smilingly ask: "What's the matter
        with you, after all, that I've had to ask you something time after time?"
        Hsi Jen opened her eyes wide. "There's nothing really the matter with me!"
        she observed; "but as you're awake, you surely had better be going over into
        the opposite room to comb your hair and wash;}
}
{
}

\Verse{44}
{
    \zA{}
}
{
    \eA{for if you dilly-dally any longer, you won't be in time." "Where shall I go
        over to?" Pao-yü inquired. Hsi Jen gave a sarcastic grin. "Do you ask me?"
        she rejoined; "do I know? you're at perfect liberty to go over wherever you
        like; from this day forth you and I must part company so as to avoid
        fighting like cocks or brawling like geese, to the amusement of third
        parties. Indeed, when you get surfeited on that side, you come over to this,
        where there are, after all, such girls as Fours and Fives (Ssu Erh and Wu
        Erh) to dance attendance upon you. But such kind of things as ourselves
        uselessly defile fine names and fine surnames." "Do you still remember this
        to-day!" Pao-yü asked with a smirk. "Hundred years hence I shall still bear
        it in mind," Hsi Jen protested;}
}
{
}

\Verse{45}
{
    \zA{}
}
{
    \eA{"I'm not like you, who treat my words as so much wind blowing by the side of
        your ears, that what I've said at night, you've forgotten early in the
        morning." Pao-yü perceiving what a seductive though angry air pervaded her
        face found it difficult to repress his feelings, and speedily taking up,
        from the side of the pillow, a hair-pin made of jade, he dashed it down
        breaking it into two exclaiming: "If I again don't listen to your words, may
        I fare like this hair-pin."}
}
{
}

\Verse{46}
{
    \zA{}
}
{
    \eA{Hsi Jen immediately picked up the hair-pin, as she remarked: "What's up with
        you at this early hour of the morning? Whether you listen or not is of no
        consequence; and is it worth while that you should behave as you do?" "How
        can you know," Pao-yü answered, "the anguish in my heart!" "Do you also know
        what anguish means?" Hsi Jen observed laughing; "if you do, then you can
        judge what the state of my heart is! But be quick and get up, and wash your
        face and be off!" As she spoke, they both got out of bed and performed their
        toilette; but after Pao-yü had gone to the drawing rooms, and at a moment
        least expected by any one, Tai-yü walked into his apartment.}
}
{
}

\Verse{47}
{
    \zA{}
}
{
    \eA{Noticing that Pao-yü was not in, she was fumbling with the books on the
        table and examining them, when, as luck would have it, she turned up the
        Chuang Tzu of the previous day. Upon perusing the passage tagged on by
        Pao-yü, she could not help feeling both incensed and amused. Nor could she
        restrain herself from taking up the pen and appending a stanza to this
        effect:  Who is that man, who of his pen, without good rhyme, made use,  A
        toilsome task to do into the Chuang-tzu text to steal,  Who for the
        knowledge he doth lack no sense of shame doth feel,  But language vile and
        foul employs third parties to abuse? At the conclusion of what she had to
        write, she too came into the drawing room;}
}
{
}

\Verse{48}
{
    \zA{}
}
{
    \eA{but after paying her respects to dowager lady Chia, she walked over to
        madame Wang's quarters. Contrary to everybody's expectations, lady Feng's
        daughter, Ta Chieh Erh, had fallen ill, and a great fuss was just going on
        as the doctor had been sent for to diagnose her ailment. "My congratulations
        to you, ladies," the doctor explained; "this young lady has fever, as she
        has small-pox; indeed it's no other complaint!" As soon as madame Wang and
        lady Feng heard the tidings, they lost no time in sending round to ascertain
        whether she was getting on all right or not, and the doctor replied: "The
        symptoms are, it is true, serious, but favourable; but though after all
        importing no danger, it's necessary to get ready the silkworms and pigs'
        tails."}
}
{
}

\Verse{49}
{
    \zA{}
}
{
    \eA{When lady Feng received this report, she, there and then, hastened to make
        the necessary preparations, and while she had the rooms swept and oblations
        offered to the goddess of small-pox, she, at the same time, transmitted
        orders to her household to avoid viands fried or roasted in fat, or other
        such heating things; and also bade P'ing Erh get ready the bedding and
        clothes for Chia Lien in a separate room, and taking pieces of deep red
        cotton material, she distributed them to the nurses, waiting-maids and all
        the servants, who were in close attendance, to cut out clothes for
        themselves.}
}
{
}

\Verse{50}
{
    \zA{}
}
{
    \eA{And having had likewise some apartments outside swept clean, she detained
        two doctors to alternately deliberate on the treatment, feel the pulse and
        administer the medicines; and for twelve days, they were not at liberty to
        return to their homes; while Chia Lien had no help but to move his quarters
        temporarily into the outer library, and lady Feng and P'ing Erh remained
        both in daily attendance upon madame Wang in her devotions to the goddess.
        Chia Lien, now that he was separated from lady Feng, soon felt disposed to
        look round for a flame.}
}
{
}

\Verse{51}
{
    \zA{}
}
{
    \eA{He had only slept alone for a couple of nights, but these nights had been so
        intensely intolerable that he had no option than to choose, for the time
        being, from among the young pages, those who were of handsome appearance,
        and bring them over to relieve his monotony. In the Jung Kuo mansion, there
        was, it happened, a cook, a most useless, good-for-nothing drunkard, whose
        name was To Kuan, in whom people recognised an infirm and a useless husband
        so that they all dubbed him with the name of To Hun Ch'ung, the stupid worm
        To.}
}
{
}

\Verse{52}
{
    \zA{}
}
{
    \eA{As the wife given to him in marriage by his father and mother was this year
        just twenty, and possessed further several traits of beauty, and was also
        naturally of a flighty and frivolous disposition, she had an extreme
        penchant for violent flirtations. But To Hun-ch'ung, on the other hand, did
        not concern himself (with her deportment), and as long as he had wine, meat
        and money he paid no heed whatever to anything. And for this reason it was
        that all the men in the two mansions of Ning and Jung had been successful in
        their attentions; and as this woman was exceptionally fascinating and
        incomparably giddy, she was generally known by all by the name To Ku Ning
        (Miss To).}
}
{
}

\Verse{53}
{
    \zA{}
}
{
    \eA{Chia Lien, now that he had his quarters outside, chafed under the pangs of
        irksome ennui, yet he too, in days gone by, had set his eyes upon this
        woman, and had for long, watered in the mouth with admiration; but as,
        inside, he feared his winsome wife, and outside, he dreaded his beloved
        lads, he had not made any advances. But this To Ku Niang had likewise a
        liking for Chia Lien, and was full of resentment at the absence of a
        favourable opportunity; but she had recently come to hear that Chia Lien had
        shifted his quarters into the outer library, and her wont was, even in the
        absence of any legitimate purpose, to go over three and four times to entice
        him on;}
}
{
}

\Verse{54}
{
    \zA{}
}
{
    \eA{but though Chia Lien was, in every respect, like a rat smitten with hunger,
        he could not dispense with holding consultation with the young friends who
        enjoyed his confidence; and as he struck a bargain with them for a large
        amount of money and silks, how could they ever not have come to terms (with
        him to speak on his behalf)? Besides, they were all old friends of this
        woman, so that, as soon as they conveyed the proposal, she willingly
        accepted it. When night came To Hun Ch'ung was lying on the couch in a state
        of drunkenness, and at the second watch, when every one was quiet, Chia Lien
        at once slipped in, and they had their assignation.}
}
{
}

\Verse{55}
{
    \zA{}
}
{
    \eA{As soon as he gazed upon her face, he lost control over his senses, and
        without even one word of ordinary greeting or commonplace remark, they
        forthwith, fervently indulged in a most endearing tête-à-tête. This woman
        possessed, who could have thought it, a strange natural charm; for, as soon
        as any one of her lovers came within any close distance of her, he speedily
        could not but notice that her very tendons and bones mollified,
        paralysed-like from feeling, so that his was the sensation of basking in a
        soft bower of love.}
}
{
}

\Verse{56}
{
    \zA{}
}
{
    \eA{What is more, her demonstrative ways and free-and-easy talk put even those
        of a born coquette to shame, with the result that while Chia Lien, at this
        time, longed to become heart and soul one with her, the woman designedly
        indulged in immodest innuendoes. "Your daughter is at home," she insinuated
        in her recumbent position, "ill with the small-pox, and prayers are being
        offered to the goddess; and your duty too should be to abstain from love
        affairs for a couple of days, but on the contrary, by flirting with me,
        you've contaminated yourself! but, you'd better be off at once from me
        here!" "You're my goddess!" gaspingly protested Chia Lien, as he gave way to
        demonstrativeness; "what do I care about any other goddess!"}
}
{
}

\Verse{57}
{
    \zA{}
}
{
    \eA{The woman began to be still more indelicate in her manner, so that Chia Lien
        could not refrain himself from making a full exhibition of his warm
        sentiments. When their tête-à-tête had come to a close, they both went on
        again to vow by the mountains and swear by the seas, and though they found
        it difficult to part company and hard to tear themselves away, they, in due
        course, became, after this occasion, mutual sworn friends.}
}
{
}

\Verse{58}
{
    \zA{}
}
{
    \eA{But by a certain day the virus in Ta Chieh's system had become exhausted,
        and the spots subsided, and at the expiry of twelve days the goddess was
        removed, and the whole household offered sacrifices to heaven, worshipped
        the ancestors, paid their vows, burnt incense, exchanged congratulations,
        and distributed presents. And these formalities observed, Chia Lien once
        more moved back into his own bedroom and was reunited with lady Feng. The
        proverb is indeed true which says: "That a new marriage is not equal to a
        long separation," for there ensued between them demonstrations of loving
        affection still more numerous than heretofore, to which we need not, of
        course, refer with any minuteness.}
}
{
}

\Verse{59}
{
    \zA{}
}
{
    \eA{The next day, at an early hour, after lady Feng had gone into the upper
        rooms, P'ing Erh set to work to put in order the clothes and bedding, which
        had been brought from outside, when, contrary to her expectation, a tress of
        hair fell out from inside the pillow-case, as she was intent upon shaking
        it. P'ing Erh understood its import, and taking at once the hair, she
        concealed it in her sleeve, and there and then came over into the room on
        this side, where she produced the hair, and smirkingly asked Chia Lien,
        "What's this?" Chia Lien, at the sight of it, lost no time in making a
        snatch with the idea of depriving her of it;}
}
{
}

\Verse{60}
{
    \zA{}
}
{
    \eA{and when P'ing Erh speedily endeavoured to run away, she was clutched by
        Chia Lien, who put her down on the stove-couch, and came up to take it from
        her hand. "You heartless fellow!" P'ing Erh laughingly exclaimed, "I conceal
        this, with every good purpose, from her knowledge, and come to ask you about
        it, and you, on the contrary, fly into a rage! But wait till she comes back,
        and I'll tell her, and we'll see what will happen." At these words, Chia
        Lien hastily forced a smile. "Dear girl!" he entreated, "give it to me, and
        I won't venture again to fly into a passion." But hardly was this remark
        finished, when they heard the voice of lady Feng penetrate into the room.}
}
{
}

\Verse{61}
{
    \zA{}
}
{
    \eA{As soon as it reached the ear of Chia Lien, he was at a loss whether it was
        better to let her go or to snatch it away, and kept on shouting, "My dear
        girl! don't let her know." P'ing Erh at once rose to her feet; but lady Feng
        had already entered the room; and she went on to bid P'ing Erh be quick and
        open a box and find a pattern for madame Wang. P'ing Erh expressed her
        obedience with alacrity; but while in search of it, lady Feng caught sight
        of Chia Lien; and suddenly remembering something, she hastened to ask P'ing
        Erh about it. "The other day," she observed, "some things were taken out,
        and have you brought them all in or not?" "I have!" P'ing Erh assented. "Is
        there anything short or not?" lady Feng inquired.}
}
{
}

\Verse{62}
{
    \zA{}
}
{
    \eA{"I've carefully looked at them," P'ing Erh added, "and haven't found even
        one single thing short." "Is there anything in excess?" lady Feng went on to
        ascertain. P'ing Erh laughed. "It's enough," she rejoined, "that there's
        nothing short; and how could there really turn out to be anything over and
        above?" "That this half month," lady Feng continued still smiling, "things
        have gone on immaculately it would be hard to vouch; for some intimate
        friend there may have been, who possibly has left something behind, in the
        shape of a ring, handkerchief or other such object, there's no saying for
        certain!"}
}
{
}

\Verse{63}
{
    \zA{}
}
{
    \eA{While these words were being spoken, Chia Lien's face turned perfectly
        sallow, and, as he stood behind lady Feng, he was intent upon gazing at
        P'ing Erh, making signs to her (that he was going) to cut her throat as a
        chicken is killed, (threatening her not to utter a sound) and entreating her
        to screen him; but P'ing Erh pretended not to notice him, and consequently
        observed smiling: "How is it that my ideas should coincide with those of
        yours, my lady; and as I suspected that there may have been something of the
        kind, I carefully searched all over, but I didn't find even so much as the
        slightest thing wrong; and if you don't believe me, my lady, you can search
        for your own self." "You fool!"}
}
{
}

\Verse{64}
{
    \zA{}
}
{
    \eA{lady Feng laughed, "had he any things of the sort, would he be likely to let
        you and I discover them!" With these words still on her lips, she took the
        patterns and went her way; whereupon P'ing Erh pointed at her nose, and
        shook her head to and fro. "In this matter," she smiled, "how much you
        should be grateful to me!" A remark which so delighted Chia Lien that his
        eyebrows distended, and his eyes smiled, and running over, he clasped her in
        his embrace, and called her promiscuously: "My darling, my pet, my own
        treasure!" "This," observed P'ing Erh, with the tress in her hand, "will be
        my source of power, during all my lifetime! if you treat me kindly, then
        well and good!}
}
{
}

\Verse{65}
{
    \zA{}
}
{
    \eA{but if you behave unkindly, then we'll at once produce this thing!" "Do put
        it away, please," Chia Lien entreated smirkingly, "and don't, on an any
        account, let her know about it!" and as he uttered these words, he noticed
        that she was off her guard, and, with a snatch, readily grabbed it adding
        laughingly: "In your hands, it would be a source of woe, so that it's better
        that I should burn it, and have done with it!" Saying this he simultaneously
        shoved it down the sides of his boot, while P'ing Erh shouted as she set her
        teeth close: "You wicked man! you cross the river and then demolish the
        bridge! but do you imagine that I'll by and by again tell lies on your
        behalf!" Chia Lien perceiving how heart-stirring her seductive charms were,
        forthwith clasped her in his arms, and begged her to be his;}
}
{
}

\Verse{66}
{
    \zA{}
}
{
    \eA{but P'ing Erh snatched her hands out of his grasp and ran away out of the
        room; which so exasperated Chia Lien that as he bent his body, he exclaimed,
        full of indignation: "What a dreadful niggardly young wench! she actually
        sets her mind to stir up people's affections with her wanton blandishments,
        and then, after all, she runs away!" "If I be wanton, it's my own look-out;"
        P'ing Erh answered, from outside the window, with a grin, "and who told you
        to arouse your affections? Do you forsooth mean to imply that my wish is to
        become your tool? And did she come to know about it would she again ever
        forgive me?" "You needn't dread her!" Chia Lien urged; "wait till my monkey
        is up, and I'll take this jealous woman, and beat her to atoms; and she'll
        then know what stuff I'm made of.}
}
{
}

\Verse{67}
{
    \zA{}
}
{
    \eA{She watches me just as she would watch a thief! and she's only to hobnob
        with men, and I'm not to say a word to any girl! and if I do say aught to a
        girl, or get anywhere near one, she must at once give way to suspicion. But
        with no regard to younger brothers or nephews, to young and old, she
        prattles and giggles with them, and doesn't entertain any fear that I may be
        jealous; but henceforward I too won't allow her to set eyes upon any man."
        "If she be jealous, there's every reason," P'ing Erh answered, "but for you
        to be jealous on her account isn't right. Her conduct is really
        straightforward, and her deportment upright, but your conduct is actuated by
        an evil heart, so much so that even I don't feel my heart at ease, not to
        say anything of her."}
}
{
}

\Verse{68}
{
    \zA{}
}
{
    \eA{"You two," continued Chia Lien, "have a mouth full of malicious breath!
        Everything the couple of you do is invariably proper, while whatever I do is
        all from an evil heart! But some time or other I shall bring you both to
        your end with my own hands!" This sentence was scarcely at an end, when lady
        Feng walked into the court. "If you're bent upon chatting," she urgently
        inquired, upon seeing P'ing Erh outside the window, "why don't you go into
        the room? and what do you mean, instead, by running out, and speaking with
        the window between?" Chia Lien from inside took up the string of the
        conversation. "You should ask her," he said. "It would verily seem as if
        there were a tiger in the room to eat her up."}
}
{
}

\Verse{69}
{
    \zA{}
}
{
    \eA{"There's not a single person in the room," P'ing Erh rejoined, "and what
        shall I stay and do with him?" "It's just the proper thing that there should
        be no one else! Isn't it?" lady Feng remarked grinning sarcastically. "Do
        these words allude to me?" P'ing Erh hastily asked, as soon as she had heard
        what she said. Lady Feng forthwith laughed. "If they don't allude to you,"
        she continued, "to whom do they?" "Don't press me to come out with some nice
        things!" P'ing Erh insinuated, and, as she spoke, she did not even raise the
        portiere (for lady Feng to enter), but straightway betook herself to the
        opposite side. Lady Feng lifted the portiere with her own hands, and walked
        into the room.}
}
{
}

\Verse{70}
{
    \zA{}
}
{
    \eA{"That girl P'ing Erh," she exclaimed, "has gone mad, and if this hussey does
        in real earnest wish to try and get the upper hand of me, it would be well
        for you to mind your skin." Chia Lien listened to her, as he kept reclining
        on the couch. "I never in the least knew," he ventured, clapping his hands
        and laughing, "that P'ing Erh was so dreadful; and I must, after all, from
        henceforth look up to her with respect!" "It's all through your humouring
        her," lady Feng rejoined; "so I'll simply settle scores with you and finish
        with it." "Ts'ui!" ejaculated Chia Lien at these words, "because you two
        can't agree, must you again make a scapegoat of me! Well then, I'll get out
        of the way of both of you!" "I'll see where you'll go and hide," lady Feng
        observed. "I've got somewhere to go!" Chia Lien added;}
}
{
}

\Verse{71}
{
    \zA{}
}
{
    \eA{and with these words, he was about to go, when lady Feng urged: "Don't be
        off! I have something to tell you." What it is, is not yet known, but,
        reader, listen to the account given in the next chapter.}
}
{
}
